% ============================================================
% config/02_styles.tex
% Styles des titres, en-tetes/pieds de page, blocs de code
% NOTE : titlesec, fancyhdr, listings, tikz et pgfornament sont
% charges dans config/00_preambule.tex, avant cet \input.
% ============================================================

% --- Bloc de chapitre (style Legrand Orange Book) ---
% CORRECTION : le mot "CHAPITRE" etait auparavant empile A L'INTERIEUR
% du bloc colore, au-dessus du chiffre, ce qui l'ecrasait. Il est
% desormais place A L'EXTERIEUR du bloc, en label vertical a gauche,
% cale sur la meme hauteur que le carre colore (cf. reference visuelle).
\newcommand{\ChapterNumberBlock}{%
  \begin{tikzpicture}[remember picture, baseline]
    % Rectangle vertical
    \node[fill=reportprimary, minimum width=2.0cm, minimum height=2.5cm, inner sep=0pt] (numbox) {};
    
    % Numéro étiré dynamiquement à 2.4cm de hauteur (remplit le bloc)
    \node[text=white, inner sep=0pt] at (numbox.center) {%
      \resizebox{!}{2.0cm}{\usefont{OT1}{pbd}{m}{n}\bfseries\thechapter}%
    };

    % Mot "CHAPITRE" étiré à la hauteur exacte du rectangle (3.0cm)
    \node[anchor=south, rotate=90, text=reportprimary, inner sep=0pt]
      at ([xshift=-4mm]numbox.west) {%
        \resizebox{2.5cm}{!}{\rmfamily\bfseries C\,H\,A\,P\,I\,T\,R\,E}%
      };
  \end{tikzpicture}%
}

\titleformat{\chapter}[display]
  {\normalfont\filleft} % Alignement de tout le bloc et du texte à droite
  {\ChapterNumberBlock} % Bloc de numéro
  {12pt}                % Espace vertical (saut de ligne) entre le bloc et le titre
  {\Huge\bfseries\color{reportprimary}} % Style du texte du titre
\titlespacing*{\chapter}{0pt}{0pt}{30pt}

\titleformat{\section}
  {\normalfont\Large\bfseries\color{reportprimary}}
  {\thesection}{1em}{}
\titleformat{\subsection}
  {\normalfont\large\bfseries}
  {\thesubsection}{1em}{}

% --- En-tete / pied de page, avec motifs ornementaux legers ---
\pagestyle{fancy}
\fancyhf{}
\fancyhead[L]{{\color{reportprimary}\pgfornament[width=\agtOrnHeaderSize]{\agtOrnHeader}}~\nouppercase{\leftmark}}
\fancyhead[R]{\rapporttitrecourt~{\color{reportprimary}\pgfornament[width=\agtOrnHeaderSize]{\agtOrnHeader}}}
\fancyfoot[C]{{\color{reportprimary}\pgfornament[width=\agtOrnFooterSize]{\agtOrnFooter}}~\thepage~{\color{reportprimary}\reflectbox{\pgfornament[width=\agtOrnFooterSize]{\agtOrnFooter}}}}
\renewcommand{\headrulewidth}{0.4pt}

% --- Mini-sommaire de chapitre (style ENSPY) ---
% Usage : appeler juste apres \chapter{...}\label{...}
\newcommand{\minisommaire}{%
  \etocsetnexttocdepth{subsection}%
  \etocsettocstyle{\section*{Contenu}}{}%
  \localtableofcontents%
  \bigskip%
}

% --- Blocs de code ---
% LIMITE ANTICIPEE : en pdfLaTeX, JetBrains Mono n'est pas disponible ;
% \ttfamily (police a chasse fixe par defaut) sert de repli fidele a
% l'esprit, sans dependance externe. En XeLaTeX/LuaLaTeX, \codefont
% (defini en 00_preambule) utilise la vraie police si installee.
\lstdefinestyle{codeagt}{
  backgroundcolor=\color{reportgraylight},
  frame=single,
  rulecolor=\color{reportbordergray},
  basicstyle=\ttfamily\small,
  breaklines=true,
  breakatwhitespace=false,
  showstringspaces=false,
  numbers=left,
  numberstyle=\tiny\color{reportgraytext},
  captionpos=b,
  tabsize=2
}
\lstset{style=codeagt}